\pdfcompresslevel=9
\NeedsTeXFormat{LaTeX2e}
\documentclass[10pt,a4paper]{article}

\usepackage[margin=0.8cm,noheadfoot,a4paper]{geometry}
\usepackage{calc}
\usepackage[english, french, italian, ngerman, spanish]{babel} % language support
\usepackage[latin1]{inputenc}
\usepackage[T1]{fontenc}
\usepackage{ae}
\usepackage{bembo}
\usepackage{ragged2e}
\usepackage{paralist}
%\usepackage{eso-pic}
\usepackage{graphicx}
\usepackage{rcsinfo}

\usepackage{hyperref}
\hypersetup {
    pdfpagemode=Full,
    colorlinks,
    urlcolor=black,
    pdftitle={GPE Flyer},
    pdfauthor={Colin Marquardt},
    pdfsubject={$Revision$, $Date$},
    pdfkeywords={GPE Palmtop Environment marketing features technology flyer},
}

% ------------------------------------------------------------------------------
% accompanying wiki page: http://handhelds.org/z/wiki/GPEMarketing

\rcsInfo $Id$
\newcommand{\flyertext}{%
\textbf{\large
  \includegraphics[height=\lineheight*3/5]{gpe-logo-simple.jpg}
GPE: The GPE Palmtop Environment
  \includegraphics[height=\lineheight*3/5]{gpe-logo-simple.jpg}
}

\smallskip
\small{
  \href{http://gpe.handhelds.org}{\texttt{http://gpe.handhelds.org}} --
  \href{mailto:gpe@handhelds.org}{\texttt{mailto:gpe@handhelds.org}} -- 
  \href{irc://freenode.net:ircd/\#gpe}{\texttt{irc://freenode.net:ircd/\#gpe}}
}

\smallskip
A Free Software GUI environment targeted at palmtop/handheld
computers, such as the HP iPAQ and the Sharp Zaurus. Think �GNOME
for handhelds�.

% GPE provides the infrastructure for developing palmtop applications,
% and includes core Personal Information Management (PIM)
% applications.

\smallskip
\textbf{Key Features}
\begin{compactitem}
   \item Custom GTK-based widget library (LGPL)
   \item Applications: calendar, to do list, address book, media
     player, notebook, games, and more (GPL)
   \item System utilities including Wi-Fi and Bluetooth manager
   \item Internationalization
   \item Display migration: �teleport� running applications to other
     displays
\end{compactitem}

\smallskip
\textbf{Technology}
\begin{compactitem}
   \item X Window System{\texttrademark}{} server with Render and
     Resize/Rotate extension (e.\,g. freedesktop.org X server, kdrive)
   \item fontconfig
   \item XFT2 with subpixel anti-aliased font support
   \item GTK+~2.2, featuring bidirectional text, display migration
     etc.
   \item freedesktop.org standards support
   \item Matchbox, a lightweight window manager designed for
     embedded devices
   \item SQLite for data storage
   \item GStreamer multimedia framework
   \item Packages available for Familiar Linux. OpenEmbedded support
     upcoming.
   \item \textbf{NEW! Synchronisation for Contacts, Todo and
     Calendar} in both directions with multisync.
\end{compactitem}

\smallskip
\textit{Users and developers welcome!}%
\hfill\tiny{Revision:~\rcsInfoRevision, \rcsInfoYear-\rcsInfoMonth-\rcsInfoDay}
}

% ------------------------------------------------------------------------------

\begin{document}

\thispagestyle{empty}
\pagestyle{empty}
\RaggedRight
\begin{table}
  \begin{tabular}{p{\textwidth/2-1cm}p{1cm}p{\textwidth/2-1cm}}%
    \flyertext{} & & \flyertext{}\\[1cm]
    \flyertext{} & & \flyertext{}
%    \\\vspace{1cm}%
  \end{tabular}
\end{table}
%% use something like this for flyers that are 3-up on landscape paper:
% \begin{table}
%   \begin{tabular}{p{\textwidth/2-1cm}p{1cm}p{\textwidth/2-1cm}p{1cm}p{\textwidth/2-1cm}}%
%     %% input another language file (that is using \renewcommand{\flyertext})
%     % \input{flyertext_es}
%     \flyertext{} & & \flyertext{} & & \flyertext{}  
%   \end{tabular}
% \end{table}

\clearpage
%% input another language file (that is using \renewcommand{\flyertext})
% This German version from Colin Marquardt is translated against
% gpe-flyer.tex Revision: 1.8

\renewcommand{\flyertext}{%
\selectlanguage{ngerman}
%
%\AddToShipoutPicture{%
%\includegraphics[width=\textwidth/2-1cm]{/opt/applications/gpe_cvs_dev/gpe/gpe-logo/gpe-logo-simple.jpg}}
\textbf{\large
  \includegraphics[height=\lineheight*3/5]{/opt/applications/gpe_cvs_dev/gpe/gpe-logo/gpe-logo-simple.jpg}
GPE: The GPE Palmtop Environment
  \includegraphics[height=\lineheight*3/5]{/opt/applications/gpe_cvs_dev/gpe/gpe-logo/gpe-logo-simple.jpg}
}

\smallskip
\small{
  \href{http://gpe.handhelds.org}{\texttt{http://gpe.handhelds.org}} --
  \href{mailto:gpe@handhelds.org}{\texttt{mailto:gpe@handhelds.org}} -- 
  \href{irc://freenode.net:ircd/\#gpe}{\texttt{irc://freenode.net:ircd/\#gpe}}
}

\smallskip

Eine GUI-Umgebung, die auf Palmtops und Handheld-Computer wie den HP
iPAQ und den Sharp Zaurus abzielt. �GNOME f�r Handhelds� ist nicht
weit weg von der Wahrheit. GPE ist Freie Software.

\smallskip
\textbf{Eigenschaften}
\begin{compactitem}
   \item Eigene GTK-basierte Widget Library (LGPL)
   \item Applikationen: Kalendar, To-do-Liste, Adre�buch,
     Medienspieler, Notizbuch, Spiele und mehr (GPL)
   \item Systemwerkzeuge wie Wi-Fi- und Bluetooth-Verwaltung
   \item Internationalisierung
   \item �Display migration�: laufende Anwendungen auf andere
  Displays verschieben
\end{compactitem}

\smallskip
\textbf{Technik}
\begin{compactitem}
   \item X Window System{\texttrademark}{}-Server mit Render- und
     Resize/Rotate-Erweiterungen (z.\,B. freedesktop.org-X-Server, kdrive)
   \item fontconfig
   \item XFT2 mit Subpixel-anti-aliasing f�r Schriften
   \item GTK+~2.2, bidirektionaler Text, �display migration�
     usw.
   \item Unterst�tzung von freedesktop.org-Standards
   \item Matchbox, ein schlanker Fenstermanager, der f�r
     Embedded-Ger�te entworfen wurde
   \item SQLite zur Datenspeicherung
   \item GStreamer Multimedia-Framework
   \item Pakete f�r Familiar Linux verf�gbar;
     OpenEmbedded-Unterst�tzung in Arbeit.
   \item \textbf{NEU! Synchronisierung f�r Kontakte, To-Do und
     Kalender} in beide Richtungen mittels multisync.
\end{compactitem}

\smallskip
\textit{Anwender und Entwickler willkommen!}%
\hfill\tiny{\rcsInfoRevision, \rcsInfoYear-\rcsInfoMonth-\rcsInfoDay}
}

%%% Local Variables: 
%%% mode: latex
%%% TeX-master: "gpe-flyer"
%%% End: 


\RaggedRight
\begin{table}
  \begin{tabular}{p{\textwidth/2-1cm}p{1cm}p{\textwidth/2-1cm}}%
    \flyertext{} & & \flyertext{}\\[1cm]
    \flyertext{} & & \flyertext{}
%    \\\vspace{1cm}%
  \end{tabular}
\end{table}

\clearpage
%% input another language file (that is using \renewcommand{\flyertext})
% this Spanish version from J.Manrique Lopez de la Fuente
% is translated against gpe-flyer.tex Revision: 1.7, Date: 2004/06/17 10:00:14

\rcsInfo $Id$
\renewcommand{\flyertext}{%
%\selectlanguage{spanish}
%
%\AddToShipoutPicture{%
%\includegraphics[width=\textwidth/2-1cm]{gpe-logo-simple.jpg}}
\textbf{\large
  \includegraphics[height=\lineheight*3/5]{gpe-logo-simple.jpg}
GPE: The GPE Palmtop Environment
  \includegraphics[height=\lineheight*3/5]{gpe-logo-simple.jpg}
}

\smallskip
\small{
  \href{http://gpe.handhelds.org}{\texttt{http://gpe.handhelds.org}} --
  \href{mailto:gpe@handhelds.org}{\texttt{mailto:gpe@handhelds.org}} -- 
  \href{irc://freenode.net:ircd/\#gpe}{\texttt{irc://freenode.net:ircd/\#gpe}}
}

\smallskip
Interface gr{\'a}fico libre orientado a ordenadores de bolsillo,
como los HP iPAQ y Sharp Zaurus. Pensado como {\guillemotright}GNOME para ordenadores
de bolsillo{\guillemotleft}.

% GPE provides the infrastructure for developing palmtop applications,
% and includes core Personal Information Management (PIM)
% applications.

% GPE provee la infraestructura necesaria para desarrollar aplicaciones
% para ordenadores de bolsillo, e incluye un n{\'u}cleo de aplicaciones 
% PIM (Gesti{\'o}n de Informaci{\'o}n Personal).

\smallskip
\textbf{Caracter{\'\i}sticas Clave}
\begin{compactitem}
   \item Librer{\'\i}a de componentes basada en GTK (LGPL)
   \item Aplicaciones: calendario, tareas, libro de direcciones, reproductor
     multimedia, notas, juegos, y m{\'a}s (GPL)
   \item Utilidades del sistema como administraci{\'o}n de Wi-Fi y Bluetooth
   \item Internacionalizaci{\'o}n
   \item Migraci{\'o}n de pantalla: {\guillemotright}teleport{\guillemotleft} ejecutando aplicaciones a otras
     pantallas
\end{compactitem}

\smallskip
\textbf{Tecnolog{\'\i}a}
\begin{compactitem}
   \item Sistema X Window {\texttrademark}{} (freedesktop.org X
     server, kdrive) con las extensiones Render y
     Redimensionar/Rotar.
   \item fontconfig.
   \item XFT2 con soporte de fuentes subpixel anti-aliased.
   \item GTK+~2.2, incluyendo soporte para texto bidireccional, migraci{\'o}n de pantalla,
     etc.
   \item Soporte de los est{\'a}ndares de freedesktop.org.
   \item Matchbox, ligero gestor de ventanas dise{\~n}ado para dispositivos
     empotrados.
   \item SQLite para el archivo de datos.
   \item GStreamer multimedia framework.
   \item Paquetes disponibles para Familiar Linux. Soporte para OpenEmbedded
     en camino.
	\item \textbf{{\textexclamdown}NOVEDAD!} Sincronizaci{\'o}n de Contactos, Tareas y Calendario en ambas direcciones con Multisync.
\end{compactitem}

\smallskip
\textit{{\textexclamdown}Usarios y desarrolladores son bienvenidos!}%
\hfill\tiny{Revision:~\rcsInfoRevision, \rcsInfoYear-\rcsInfoMonth-\rcsInfoDay}
}

%%% Local Variables: 
%%% mode: latex
%%% TeX-master: "gpe-flyer"
%%% End: 


\RaggedRight
\begin{table}
  \begin{tabular}{p{\textwidth/2-1cm}p{1cm}p{\textwidth/2-1cm}}%
    \flyertext{} & & \flyertext{}\\[1cm]
    \flyertext{} & & \flyertext{}
%    \\\vspace{1cm}%
  \end{tabular}
\end{table}

\clearpage
%% input another language file (that is using \renewcommand{\flyertext})
% this French version from Michael Opdenacker
% is translated against gpe-flyer.tex Revision: 1.10, Tue Apr 13 19:39:44 2004 UTC
\rcsInfo $Id$
\renewcommand{\flyertext}{%
%\selectlanguage{french}
%
%\AddToShipoutPicture{%
%\includegraphics[width=\textwidth/2-1cm]{gpe-logo-simple.jpg}}
\textbf{\large
  \includegraphics[height=\lineheight*3/5]{gpe-logo-simple.jpg}
GPE: The GPE Palmtop Environment
  \includegraphics[height=\lineheight*3/5]{gpe-logo-simple.jpg}
}

\smallskip
\small{
  \href{http://gpe.handhelds.org}{\texttt{http://gpe.handhelds.org}} --
  \href{mailto:gpe@handhelds.org}{\texttt{mailto:gpe@handhelds.org}} -- 
  \href{irc://freenode.net:ircd/\#gpe}{\texttt{irc://freenode.net:ircd/\#gpe}}
}

\smallskip
Un Logiciel Libre d'environment graphique destin{\'e} aux ordinateurs de
poche, tels que les iPAQ d'HP et les Zaurus de Sharp.
Consid{\'e}rez le comme {\guillemotright}GNOME pour ordinateurs de poche{\guillemotleft}

% GPE provides the infrastructure for developing palmtop applications,
% and includes core Personal Information Management (PIM)
% applications.

\smallskip
\textbf{Caract{\'e}ristiques principales}
\begin{compactitem}
   \item Biblioth{\`e}que d{\'e}di{\'e}e de composants graphiques GTK (LGPL)
   \item Applications: agenda, gestionnaire de t{\^a}ches, carnet d'adresses,
     lecteur multim{\'e}dia, bloc-notes, jeux et plus encore (GPL)
   \item Outils d'administration syst{\`e}me, dont le gestionnaire Wi-Fi et Bluetooth
   \item Internationalisation
   \item Migration d'affichage: d{\'e}porter des applications en cours d'ex{\'e}custion vers d'autres {\'e}crans.
\end{compactitem}

\smallskip
\textbf{Technologie}
\begin{compactitem}
   \item Serveur X Window System{\texttrademark}{} (freedesktop.org X
     server, kdrive) avec les extensions Render et
     Retailler/Tourner
   \item fontconfig, XFT2 avec support de polices liss{\'e}es
   \item GTK+~2.2, incluant la gestion de texte bidirectionnel, le d{\'e}port d'{\'e}cran, etc.
   \item Prise en compte des standarts de freedesktop.org
   \item Matchbox, un gestionnaire de fen{\^e}tres l{\'e}ger con\c{c}u pour les syst{\`e}mes embarqu{\'e}s.
   \item SQLite pour le stockage de donn{\'e}es
   \item Bo{\^i}te {\`a} outils multim{\'e}dia GStreamer
   \item Paquetages disponibles pour la distribution Familiar Linux. Bient{\^o}t support d'OpenEmbedded.
   \item \textbf{NOUVEAU! Synchronisation du carnet d'adresses, du gestionnaire de t{\^a}ches et de l'agenda} dans
   	les deux sens avec multisync.
\end{compactitem}

\smallskip
\textit{{\textexclamdown}Utilisateurs et d{\'e}veloppeurs bienvenus!}%
\hfill\tiny{Revision:~\rcsInfoRevision, \rcsInfoYear-\rcsInfoMonth-\rcsInfoDay}
}

%%% Local Variables: 
%%% mode: latex
%%% TeX-master: "gpe-flyer"
%%% End: 


\RaggedRight
\begin{table}
  \begin{tabular}{p{\textwidth/2-1cm}p{1cm}p{\textwidth/2-1cm}}%
    \flyertext{} & & \flyertext{}\\[1cm]
    \flyertext{} & & \flyertext{}
%    \\\vspace{1cm}%
  \end{tabular}
\end{table}

\clearpage
%% input another language file (that is using \renewcommand{\flyertext})
% Italian version by Giansalvo Gusinu <giansalvo at gusinu.net>
% translation of gpe-flyer.tex Revision: 1.10, Tue Apr 13 19:39:44 2004 UTC
\rcsInfo $Id$
\renewcommand{\flyertext}{%
%\selectlanguage{french}
%
%\AddToShipoutPicture{%
%\includegraphics[width=\textwidth/2-1cm]{gpe-logo-simple.jpg}}
\textbf{\large
  \includegraphics[height=\lineheight*3/5]{gpe-logo-simple.jpg}
GPE: Ambiente grafico per palmari
  \includegraphics[height=\lineheight*3/5]{gpe-logo-simple.jpg}
}

\smallskip
\small{
  \href{http://gpe.handhelds.org}{\texttt{http://gpe.handhelds.org}} --
  \href{mailto:gpe@handhelds.org}{\texttt{mailto:gpe@handhelds.org}} -- 
  \href{irc://freenode.net:ircd/\#gpe}{\texttt{irc://freenode.net:ircd/\#gpe}}
}

\smallskip
Ambiente grafico destinato a computer palmari tipo iPAQ di HP e Zaurus di Sharp.
Rilasciato con licenza 'free/open source'. Pensate a {\guillemotright}GNOME per palmari'{\guillemotleft}.

\smallskip
\textbf{Caratteristiche principali}
\begin{compactitem}
   \item Libreria specifica di widget grafici basata su GTK (LGPL)
   \item Applicazioni: calendario, agenda, gestione di task, rubrica indirizzi,
     lettore multimediale, bloc notes, giochi e altro (GPL)
   \item Utilit{\`a} di sistema tra cui WiFi manager et Bluetooth manager
   \item Internazionalizzazione
   \item Migrazione del display: trasferimento dell'output di un programma 
            in esecuzione su un altro schermo.
\end{compactitem}

\smallskip
\textbf{Tecnologia}
\begin{compactitem}
   \item Server X Window System{\texttrademark}{} (freedesktop.org X
     server, kdrive) con estensioni Render et Resize/Rotate
   \item fontconfig, XFT2 con supporto per font con subpixel anti-aliased
   \item GTK+~2.2, supporto per testo bidirezionale, migrazione del display, etc.
   \item supporto degli standard freedesktop.org
   \item Matchbox, un window manager leggero progettato per apparecchi embedded.
   \item SQLite per database
   \item framework multimediale GStreamer
   \item Pacchetti disponibili per la distribuzione Linux Familiar. Supporto per OpenEmbedded in preparazione.
   \item \textbf{NUOVO! Sincronizzazione dei Contatti, TODO e Calendario} 
   bidirezionale con multisync.
\end{compactitem}

\smallskip
\textit{{\textexclamdown}Utenti e sviluppatori sono benvenuti!}%
\hfill\tiny{Revision:~\rcsInfoRevision, \rcsInfoYear-\rcsInfoMonth-\rcsInfoDay}
}

%%% Local Variables: 
%%% mode: latex
%%% TeX-master: "gpe-flyer"
%%% End: 


\RaggedRight
\begin{table}
  \begin{tabular}{p{\textwidth/2-1cm}p{1cm}p{\textwidth/2-1cm}}%
    \flyertext{} & & \flyertext{}\\[1cm]
    \flyertext{} & & \flyertext{}
%    \\\vspace{1cm}%
  \end{tabular}
\end{table}

\end{document}

% From: Andreas Matthias <amat@KABSI.AT>
% Subject: Re: Subject: 4 gleiche Flyer auf ein A4-Blatt?
% To: TEX-D-L@LISTSERV.DFN.DE
% User-Agent: Gnus/5.09 (Gnus v5.9.0) Emacs/21.3.50
% X-Gnus-Mail-Source: file:/var/spool/mail/colin
% Message-ID:  <m3u1ekkpjz.fsf@h081217002064.dyn.cm.kabsi.at>
% 
% Colin Marquardt wrote:
% 
% > Herbert Voss <Herbert.Voss@ALUMNI.TU-BERLIN.DE> writes:
% > 
% >> Colin Marquardt schrieb:
% >> 
% >>> Ich muss einen Flyer setzen
% >>> (<http://handhelds.org/z/wiki/GPEMarketing>, wen es interessiert).
% >> 
% >> z.B. in eine Tabelle setzen
% > 
% > So habe ich das letztendlich auch gemacht.
% 
% 
% Ich hatte eigentlich gehofft, dass sich irgendwo noch ein
% Tool f�r diese Aufgabe findet, weshalb ich meine primitive
% Bastlerl�sung nicht posten wollte. Sie ist aber vielleicht
% doch etwas einfacher und allgemeiner als alles mit Tabellen
% zu versuchen. Daher hier doch kurz ein Beispiel, wie man es
% mit pdfpages.sty in 3 Schritten zusammenpfuschen kann.
% 
% Ciao
% Andreas
% 
% 
% %%% %%% dummy.tex %%% %%%
% \documentclass{article}
% \usepackage[a6paper]{geometry}
% \begin{document}
% First page \ldots \vfill First page \ldots\newpage
% Second page \ldots \vfill Second page \ldots
% \end{document}
% %%% %%% dummy.tex %%% %%%
% 
% 
% %%% %%% dummy_2.tex %%% %%%
% \documentclass[a4paper]{article}
% \usepackage{pdfpages}
% \begin{document}
% \includepdf[pages=-, doublepages, nup=1x2, landscape]{dummy.pdf}
% \end{document}
% %%% %%% dummy_2.tex %%% %%%
% 
% 
% %%% %%% dummy_4.tex %%% %%%
% \documentclass[a4paper]{article}
% \usepackage{pdfpages}
% \begin{document}
% \includepdf[pages=-, doublepages, nup=1x2]{dummy_2.pdf}
% \end{document}
% %%% %%% dummy_4.tex %%% %%%
